% Independent continuation after B.2.26.  Retains the source A4 map separately.
\subsection{A harmonic character detected in the full descent kernel}

Keep $B_0=C_2(s)*C_3(r)$, $F_0$, and the discrete sixfold
kernel $K_0$ of B.2.15.  Thus $U_0=B_0/K_0$, and its
completion is the actual profinite quotient $U=B/K$.
Write $q:B_0\twoheadrightarrow A_4$ for the specified
source modular quotient
\[
 q(s)=(12)(34),\qquad q(r)=(123).
 \tag{B.Z1}
\]
This is a source quotient, retained in addition to the
six target-coordinate quotients.  Put
\[
 K_0^+=K_0\cap\ker q,\quad U_0^+=B_0/K_0^+,
 \qquad K^+=K\cap\ker q,\quad U^+=B/K^+.
 \tag{B.Z2}
\]
The maps in (B.Z2) are induced by the same named $s,r$.

\begin{result}{B.2.27}{Proved}{Two independent characters on the actual descent kernel}
There are actual elements $a,k\in H\cap\ker D_P$ such that
\[
 \begin{array}{c|cc}
       &\Delta&\overline d_2\\\hline
 a     &1&1\\
 k     &1&0\\
 ak    &0&1.
 \end{array}
 \tag{B.Z3}
\]
In particular $\overline d_2$ does not belong to the
subspace spanned by $\Delta D_P^n$ for $n\geq0$.
\end{result}
\begin{proof}
We first identify the whole joint quotient in (B.Z2).
Use the explicit words
\[
 c=x^{-2}y^{-2}x^2y^2,\qquad
 w=c^{-1}z^{-2}cz^2,\qquad z=(xy)^{-1}.
 \tag{B.Z4}
\]
Every one of the six source-deck substitutions sends
one of the three peripheral classes $x,y,z$, after
the target map $\varphi_P$, to a conjugate of the
order-two target element $s$.  The corresponding
squared input in (B.Z4) is one.  Each of the three
possibilities makes the displayed iterated commutator
one.  Hence $w\in K_0\cap F_0'$ by literal word
evaluation in all six full target groups.

The source quotient (B.Z1), evaluated at
$x=(sr)^{-2}$ and $y=sxs$, instead gives
\[
                         q(w)=q(s)=(12)(34).
 \tag{B.Z5}
\]
This four-letter permutation calculation is reproduced
independently in \path{check_source_affine_brunnian.py}.
The script also checks the three choices obtained by
rotating this nonzero element by $r$.  Since $K_0$ is
normal, (B.Z5) implies $q(K_0)=V_4$.  The opposite
containment follows from $K_0\subseteq B_0'$ and
$A_4'=V_4$.  Consequently the actual joint image is
the whole fibre product over $A_4/V_4=C_3$:
\[
 U_0^+=U_0\mathbin{\times}_{C_3}A_4
            \cong V_4\rtimes U_0.
 \tag{B.Z6}
\]
Here $U_0\to C_3$ is the abelian character taking
$s\mapsto0$ and $r\mapsto1$.  The action on $V_4$
is its usual faithful order-three action.  Choosing
the section of $A_4\to C_3$ generated by the named
$q(r)$ makes (B.Z6) explicit: $(v,u)$ corresponds
to $(u,vq(r)^{\chi_3(u)})$.  Thus the actual source
elements and the element represented by (B.Z4) are
\[
 s^+=(v_s,s_U),\qquad r^+=(0,r_U),\qquad
 b=(v_s,1),\qquad v_s=q(s)\ne0.
 \tag{B.Z7}
\]
In particular $b\in C_{J_0^+}(s^+)$, where
$J_0^+=F_0/K_0^+$, and $b$ maps to one in $U_0$.

We need the full central lift of this particular
centralizer element.  The required integral central
kernel can be controlled without substituting a
finite quotient for $U_0^+$.

First, the discrete central kernel
\[
 M_0=K_0/[B_0,K_0]\cong H_2(U_0,\mathbb Z)
 \tag{B.Z8}
\]
is finite by B.2.15.  It is generated by the images
of the two discrete torsion correction maps
$\kappa_{s,0},\kappa_{r,0}$.  We spell out the
discrete scope of the earlier proof.  The ternary
calculation (B.S60) makes $(M_0)_{(3)}$ cyclic.
The discrete centralizer surjection in the proof
of B.2.20 gives an order-three subgroup mapping
onto the order-three kernel at $E$; the cyclic-group
argument following (B.S60) therefore gives
$(M_0)_{(3)}=\operatorname{im}\kappa_{r,0}\cong C_3$.
For the binary part, take a character of
$M_0/(\operatorname{im}\kappa_{s,0}
\operatorname{im}\kappa_{r,0})$ to $\mathbb F_2$.
Its pushout defines exactly the marked class in
$H^2(U_0,\mathbb F_2)$ used in (B.S66).
The diagonal centralizers (B.S68)--(B.S71) are
discrete words, and kill its restriction to $L^6$.
The actual filtration isomorphism (B.S73) and the
rank-four certificate (B.S74) apply to that class.
The centralizer surjection (B.S75) is already stated
for $J_0$, and is proved by solving the three
coordinate pairs in the discrete group $P^6$.
Thus the resulting class is zero.  The injective
discrete transgression (B.S76) makes the character
zero on $M_0$ directly, without a completion step.
The finite reduced two-group is therefore trivial.
This proves the asserted generation of $M_0$.
Since the two commutator maps take values of orders
dividing two and three respectively, we obtain
\[
             2(M_0)_{(2)}=0,
                 \qquad 3(M_0)_{(3)}=0.
 \tag{B.Z9}
\]

Next regard $V_4$ as its two-dimensional
$\mathbb F_2[U_0]$ module via $\chi_3$.  One has
\[
                         H^1(U_0,V_4)=0.
 \tag{B.Z10}
\]
Indeed inflation along $B_0\twoheadrightarrow U_0$
is injective on this first cohomology.  On $B_0$,
the generator $s$ acts trivially, and $r-1$ is
invertible on $V_4$.  Every cocycle has a unique
representative modulo coboundaries with value zero
at $r$; its value at $s$ can be any element of
$V_4$.  Every nonzero choice defines the actual
map to $V_4\rtimes C_3=A_4$, with the named
$r$ fixed and $s$ taken to that nonzero element.
The three such choices are related by the
$C_3$-equivariant scalar automorphisms of
$V_4=\mathbb F_4$.  By (B.Z5), each takes $w$
to a nonzero element.  None can factor through
$U_0$, because $w\in K_0$.  This proves (B.Z10).
The module is self-dual by its invariant nondegenerate
alternating pairing.  Duality over $\mathbb F_2$
therefore also gives $H_1(U_0,V_4)=0$.
Furthermore $(V_4)_{U_0}=0$ because $r-1$ is
invertible.

Apply the integral homology spectral sequence to
the actual split extension (B.Z6).  In total degree
two, its three terms are
$H_2(U_0,\mathbb Z)$, $H_1(U_0,V_4)=0$, and
$H_0(U_0,H_2(V_4,\mathbb Z))=C_2$.
The first term has no nonzero outgoing differential:
the only possible one lands in $(V_4)_{U_0}=0$.
The final $C_2$ survives as well.  To see this with
actual maps, the projection $U_0^+\to A_4$ has
the section
\[
                  (v,q(r)^j)\longmapsto(v,r_U^j).
 \tag{B.Z11}
\]
The nonzero element of $H_2(V_4,\mathbb Z)$ is
detected in $A_4$ by the central extension
$\operatorname{SL}_2(\mathbb F_3)\to A_4$:
its restriction to $V_4$ is $Q_8\to V_4$, and
the lifts of two independent elements have
nontrivial central commutator.  Thus no incoming
differential kills that element.  The remaining
extension of the two terms splits through the
section $U_0\to U_0^+$, $u\mapsto(0,u)$.
We have proved the actual isomorphism
\[
 H_2(U_0^+,\mathbb Z)
       \cong H_2(U_0,\mathbb Z)\oplus C_2.
 \tag{B.Z12}
\]
Both projections and sections have just been specified.
Since the $U_0$-coinvariants of $V_4$ are zero,
$(U_0^+)^{\rm ab}=U_0^{\rm ab}=C_6$.
The same five-term sequence as in (B.Z8) identifies
\[
 M_0^+=K_0^+/[B_0,K_0^+]
          \cong H_2(U_0^+,\mathbb Z).
 \tag{B.Z13}
\]
In particular (B.Z9)--(B.Z12) show that its entire
binary part is elementary abelian.

We now compute $\kappa_{s,0}^+(b)$ in this full
central kernel.  With $\mathbb F_2$ coefficients,
inflation from $U_0$ is injective by the section,
and (B.Z12) together with the universal-coefficient
sequence shows that its cokernel in
$H^2(U_0^+,\mathbb F_2)$ has dimension one.
A representative of the extra class is the pullback
of $\operatorname{SL}_2(\mathbb F_3)\to A_4$.
For an inflated class, $b$ projects to one in
$U_0$, so the commutator of its lift with a lift
of $s^+$ is one.  For the additional class,
$q(b)=q(s^+)=q(s)$; their lifts in
$\operatorname{SL}_2(\mathbb F_3)$ are powers of
one element of order four and also commute.
This proves zero evaluation on
$\kappa_{s,0}^+(b)$ for every binary central class.

The marked-source condition causes no omitted class
in this argument.  The quaternion class restricts
nontrivially to the named source $C_2$, because
the indicated lift has order four.  Adding the
pullback of the class $C_4\to C_2$ along the
abelian quotient $U_0\to C_2$ corrects that
restriction.  This added class has zero commutator
on $b$, which projects to one in $U_0$ and has
trivial abelianization.  Hence every canonical
transgression class has the same zero evaluation.
The element $\kappa_{s,0}^+(b)$ has order at most
two.  Since the binary part in (B.Z13) is
elementary, its binary characters separate it.
Consequently
\[
                         \kappa_{s,0}^+(b)=1.
 \tag{B.Z14}
\]

This identity transfers through a direct source
morphism to the actual profinite central cover.
Indeed (B.Z14), with $w$ as its specified lift,
states the literal containment
$[w,s]\in[B_0,K_0^+]$.
The natural map $B_0\to B$ sends $K_0^+$ into
$K^+$, so it sends this commutator into
$\overline{[B,K^+]}$.  It therefore defines the
required centralizer identity in
\[
              (U^+)^{\mathrm c}=B/\overline{[B,K^+]}
 \tag{B.Z15}
\]
without a further comparison of completion kernels.
The word $w$ of (B.Z4) supplies a pure lift
$\widetilde b$ of $b$ to (B.Z15), and
$[\widetilde b,s]=1$.
The profinite joint quotient itself is
$U^+=U\times_{C_3}A_4$: the same word gives
$q(K)=V_4$, exactly as in (B.Z5)--(B.Z6).

Take the identity automorphism of (B.Z15), with
the positive $s$-frame $\widetilde b$, the
positive $r$-frame one, and cusp frame one.
B.2.14 supplies a full positive cusp-fixed
automorphism $\Phi$ of $B$ with precisely these
$U^+$-action and frame values.  Write
$f=s\Phi(s)$, and write its pure frames as
$\beta,A$.  Before commutator normalization one has
\[
 \Phi_U=1,\quad f,\beta,A\in K,\qquad
                   q(\beta)=q(s),\quad q(A)=1.
 \tag{B.Z16}
\]
These are assertions in the full quotient $U$,
obtained from the full central-cover image equality.

We retain the normalization explicitly.  Put
$\beta^{\rm ab}=(u,v)$ and $t=v-u$.
The first two of the six maps take the pure
abelian coordinates to $3u+v$ and $u+3v$ in
the target $B^{\rm ab}=C_6$.  As $\beta\in K$,
both vanish.  They imply $u,v\equiv0\pmod3$
and $u\equiv v\pmod2$, hence
$t\in6\widehat{\mathbb Z}$.
The corrected automorphism
$\Phi_0=\operatorname{Inn}(x^t)\Phi$
has its unique normalized gentle parameter
$a_0=(1,f_0)\in H$.  Its descended action changes
from the identity by
$\operatorname{Inn}(\varphi_P(x)^t)
=\operatorname{Inn}(s^t)=1$.
Thus $D_P(a_0)=1$ literally.  Also $q(x)$ has
order three, so this normalization preserves
$q(\beta)=q(s)$ and $q(A)=1$ in (B.Z16).

In the projective level-three representation,
$q(A_{a_0})=1$ forces the correcting shear
$q(V_{a_0})$ to be one: it belongs to the
order-three shear subgroup and centralizes the
distinct named cubic subgroup, whose intersection
with that shear subgroup is trivial.  Therefore
$q(W_{a_0})=q(s)$ and
$\epsilon_2(W_{a_0})=1$.
The exact descent identity (B.C24) gives
$\Psi_{3,2}(a_0)/2=
\epsilon_2(W_{D_P(a_0)})=0$ modulo two.
Hence $\overline d_2(a_0)=d_2(a_0)=1$.

The existing element $k$ of B.1.3 is in
$H\cap\ker D_P$ and has
$(\Delta(k),\overline d_2(k))=(1,0)$ by B.2.26.
Set $a=a_0 k^{1+\Delta(a_0)}$, where the exponent
is the element zero or one of $\mathbb F_2$.
This gives precisely (B.Z3), while preserving
membership in the full descent kernel.

Finally $\chi_n(ak)=0$ for every $n\geq0$:
the case $n=0$ is $\Delta(ak)=0$, and all later
cases follow from $D_P(ak)=1$.
But $\overline d_2(ak)=1$.  Evaluation on this
actual full profinite element proves the last
assertion of the result.
\end{proof}
